\section{Checks after acquisition and on synthetic logs}
\label{app:checks}

% Moved from Sections 4.2, 4.3 and 5.1 of the main body on 24 October 2026, without changing the procedure. The text
% is that of the main body before the move, except for the references back to it.

\subsection{Checks after acquisition}
\label{sec:data:checks}

Before any value of matching or utility was computed, the preregistration allowed only a fixed list of checks: the
counts of files, rows, writers and excluded writers; the column names and values; the share of logs whose final draft,
rebuilt from the unprocessed log, matches the distributed text; the error of the segment rule
(\S\ref{sec:attackers:segments}); the structure of ENJA15; the inference time of KeyTAR on one log; and a verification
of our formulas for the process measures. The checks report counts only; the scripts print no writer identifier, no
text and no log.

\paragraph{Escaped characters.}
The distributed logs write four characters (newline, tab, double quote and backslash) in the text change and key
columns as two-character escape sequences, which the procedure did not anticipate. Read as they are, a newline is
inserted as two characters and the rebuilt draft drifts from the cursor position. We undo the escapes when the logs are
read and change nothing else in the procedure; this is reported as a deviation, with the values of the checks before
and after the change. After it, the final draft rebuilt from the unprocessed log matches the distributed text exactly,
after normalisation of whitespace, for 4,842 of the 4,992 writers (97.0\%), against 72 before. The other 150 writers
stay in full-text matching, and accuracy is also reported separately for writers whose rebuilt draft matches and
writers whose draft does not.

\paragraph{Other checks.}
The action time of modifier keys, which the corpus description gives as zero~\cite{tian2025klicke}, is zero for Shift,
Control and Alt but not for most CapsLock and Meta presses. No attacker uses hold times and the coarser forms drop
them, so the procedure is unchanged. The corpus does not distribute the process measures computed by Inputlog, so the
verification of our formulas against them could not be done (\S\ref{sec:protocol:utility}). For KeyTAR, the default
settings of the artifact were recorded, and the inference on one log fitted in the memory of a 16\,GB GPU.

\paragraph{Cases that the segment rule misses.}
The rule ``the next press after Control'' misses shortcuts made with several keys while Control is held,
such as two successive undos, and pastes made from a menu: in 566 of the 1,133 paste events of the corpus, the
preceding press, ignoring Shift, Alt and CapsLock, is not Control.

\subsection{A check on synthetic logs before registration}
\label{app:checks:synthetic}

Before the preregistration was fixed, we checked on synthetic logs, with no human data, whether the lengths of words
can be recovered from columns other than the key categories. The synthetic logs included deletion, deletion of a
selection, replacement, paste, cursor movement by arrows, Home, End and mouse, and undo. In pairs of synthetic logs
that share the category sequence and the press times and differ only in hidden choices, such as the target of a mouse
click or the pasted words, forms R3, R3-u, R3-t and R7 gave identical output, while R0, R1 and R2 differed in every
pair (R2 through the position and the word count). The release times were at first kept in these forms; in a
simulation in which a drag's duration appears in its hold time, the output differed through them, and they were then
dropped. A drag's duration, and that of a held key when the next press follows its release, still remains as a lower
bound of the next press interval, which a form that releases press times cannot remove. In R2, rebuilding the draft
from the cursor position recovered the exact sequence of word lengths in 64\% of the synthetic logs with light
editing. These checks led to the present definitions of the forms, the categories and the segment rule, which are part
of the registered plan.
